\documentclass[10pt,letterpaper,twocolumn]{article}

\usepackage[T1]{fontenc}
\usepackage[utf8]{inputenc}
\usepackage[margin=0.55in,columnsep=0.20in]{geometry}
\usepackage{amsmath,amssymb,mathtools}
\usepackage[dvipsnames]{xcolor}
\usepackage{booktabs,tabularx}
\usepackage{enumitem}
\usepackage{graphicx}
\usepackage{wrapfig}
\usepackage{microtype}
\usepackage[hidelinks]{hyperref}

\setlength{\parindent}{0pt}
\setlength{\parskip}{1.5pt}
\setlength{\abovedisplayskip}{2pt plus 0.5pt minus 0.5pt}
\setlength{\belowdisplayskip}{2pt plus 0.5pt minus 0.5pt}
\setlength{\abovedisplayshortskip}{1pt plus 0.5pt}
\setlength{\belowdisplayshortskip}{1.5pt plus 0.5pt minus 0.5pt}
\setlength{\jot}{2pt}
\setlength{\intextsep}{2pt plus 1pt minus 1pt}
\setlist{nosep,leftmargin=*,topsep=1pt,partopsep=0pt}
\definecolor{peblue}{HTML}{E6F0FA}
\definecolor{peamber}{HTML}{FFF2CC}
\definecolor{peviolet}{HTML}{EDE5F6}
\definecolor{pegreen}{HTML}{E2F2E6}
\newcommand{\pehead}[1]{\textbf{#1.}\hspace{0.30em}}
\newcommand{\pemathmark}[2]{%
  \begingroup\setlength{\fboxsep}{0.55pt}%
  \mathchoice
    {\colorbox{#1}{$\displaystyle\color{black}\mathstrut #2$}}%
    {\colorbox{#1}{$\textstyle\color{black}\mathstrut #2$}}%
    {\colorbox{#1}{$\scriptstyle\color{black}\mathstrut #2$}}%
    {\colorbox{#1}{$\scriptscriptstyle\color{black}\mathstrut #2$}}%
  \endgroup
}
\newcommand{\petextmark}[2]{%
  \begingroup\setlength{\fboxsep}{0.55pt}%
  \colorbox{#1}{\textcolor{black}{\strut #2}}%
  \endgroup
}
\pagestyle{empty}
\frenchspacing
\raggedbottom
\emergencystretch=1em

\begin{document}

{\fontsize{9.5}{10.5}\selectfont\bfseries \href{https://arxiv.org/abs/1706.03762v7}{Attention Is All You Need}%
\hfill{\normalfont Vaswani et al. (2017)}\par}
\vspace{1pt}

\pehead{Problem \& contribution} Recurrent sequence-to-sequence models propagate state one position at a time, which serializes training within each example and makes distant tokens interact through long computational paths. The paper introduces the \emph{Transformer}: an encoder--decoder model that uses attention for all cross-position communication and applies the same small feed-forward transformation independently at every position. It therefore processes all training positions in parallel while letting any position directly select information from any other; the decoder remains autoregressive at generation time.

\pehead{Routing mechanism} For query vectors $Q$, key vectors $K$, values $V$, and key width $d_k$, scaled dot-product attention is
\[
 \operatorname{Attn}(Q,K,V)=
 \pemathmark{pegreen}{\operatorname{softmax}}
 \!\left(\pemathmark{peblue}{QK^{\mathsf T}}
 \pemathmark{peamber}{/\sqrt{d_k}}\right)
 \pemathmark{peviolet}{V}.
\]
The \petextmark{peblue}{query--key scores} measure which positions match; the \petextmark{peamber}{scale factor} keeps their variance near one when components have unit variance, preventing large logits from saturating softmax; \petextmark{pegreen}{row-wise normalization} turns each query row into weights; and the weighted \petextmark{peviolet}{value mixture} retrieves the selected information. Thus each head acts as differentiable content-addressed memory: queries are lookups, keys are addresses, and values are payloads. Encoder self-attention derives $Q,K,V$ from the same sequence; masked decoder self-attention does likewise but blocks future keys; cross-attention derives $Q$ from the decoder and $K,V$ from encoder outputs. One mixture can blur different relations, so $h$ heads learn separate projections:
\[
\begin{aligned}
 H_i&=\operatorname{Attn}(QW_i^Q,KW_i^K,VW_i^V),\\
 \operatorname{MHA}&=\operatorname{Concat}_{i=1}^{h}(H_i)W^O.
\end{aligned}
\]
Each head uses reduced dimensions, so several routing patterns cost roughly one full-width head rather than $h$ full models.

\pehead{Architecture} Each of six encoder layers applies self-attention then a position-wise feed-forward network. Each of six decoder layers applies masked self-attention, encoder--decoder attention, then the feed-forward network. The decoder receives the previously emitted symbols shifted right; a triangular mask assigns future-token logits $-\infty$, preserving the autoregressive factorization, while cross-attention lets each generated position query the encoded input. A tied embedding/output matrix maps final decoder states to vocabulary logits. Every sublayer uses the residual form $\operatorname{LayerNorm}(x+\operatorname{Sublayer}(x))$. The feed-forward map $\max(0,xW_1+b_1)W_2+b_2$ is shared across positions but not layers. Because attention itself is order-agnostic, fixed sinusoidal coordinates are added to embeddings (scaled by $\sqrt{d_{\rm model}}$): $PE_{p,2i}=\sin(p/10000^{2i/d})$ and $PE_{p,2i+1}=\cos(p/10000^{2i/d})$. A fixed offset can be represented as a linear transformation of these codes, giving the model a usable representation of relative displacement; learned positions performed almost identically in the reported ablation. The base model uses $d_{\rm model}=512$, $d_{\rm ff}=2048$, $h=8$, and $d_k=d_v=64$. Training minimizes teacher-forced next-token cross-entropy with label smoothing $0.1$. Adam uses $\beta_1=0.9$, $\beta_2=0.98$, and $\epsilon=10^{-9}$ with learning rate $d_{\rm model}^{-1/2}\min(t^{-1/2},t\,4000^{-3/2})$: linear warmup stabilizes early optimization, then inverse-square-root decay lowers the step size. Residual and attention dropout are $0.1$ in the base model.

Its computational tradeoff helps explain the design:
\begin{center}
\begin{tabularx}{\columnwidth}{@{}lXXX@{}}
\toprule
Layer & Work & Sequential & Max path\\
\midrule
Self-attention & $O(n^2d)$ & $O(1)$ & $O(1)$\\
Recurrent & $O(nd^2)$ & $O(n)$ & $O(n)$\\
Convolution & $O(knd^2)$ & $O(1)$ & $O(\log_k n)$\\
\bottomrule
\end{tabularx}
\end{center}
Dense attention was cheaper than recurrence when $n<d$ in the translation setting and exposed short gradient paths, but its quadratic sequence term is the central scaling weakness. Multi-head projections preserve this routing advantage while allowing different representation subspaces and relative positions to be selected simultaneously.

\pehead{Evidence \& scope} On WMT 2014 English--German (about 4.5M sentence pairs), the 65M-parameter base model reports 27.3 BLEU at an estimated $3.3\!\times\!10^{18}$ training FLOPs after 100k steps (12 hours on eight P100 GPUs); the 213M-parameter big model reports 28.4 BLEU, over 2 BLEU above prior reported systems including ensembles, after 300k steps (3.5 days). The big variant also doubles model width to 1024, doubles heads to 16, doubles feed-forward width to 4096, and raises dropout to $0.3$, so its gain does not isolate one change. For English--French (36M pairs), the abstract and table report 41.8 BLEU while the prose says 41.0, an internal source inconsistency. The reported translations average recent checkpoints and use a shared byte-pair vocabulary, beam size 4, and length penalty $0.6$; these conditions demonstrate strong quality/compute tradeoffs in these tasks, not architecture-independent superiority. Development ablations support multiple heads without showing that more is always better: one head scored 24.9 BLEU, 8/16 heads 25.8, and 32 heads 25.4. Label smoothing worsened perplexity but improved accuracy and BLEU, warning that likelihood and task metrics need not rank models identically. A four-layer Transformer reached 91.3 F1 on WSJ-only constituency parsing and 92.7 with semi-supervision, evidence of transfer beyond translation rather than proof of universal superiority.

\pehead{Limits} Dense attention trades parallel training and short paths for $O(n^2)$ memory and work; the paper leaves local attention for long inputs to future work. Autoregressive decoding still generates tokens sequentially, and the experiments do not report matched end-to-end inference latency or memory. Sinusoidal extrapolation beyond trained lengths is a motivation, not an established result. Attention plots show syntactic-looking and long-range patterns, but qualitative examples do not prove that attention weights faithfully explain predictions. The results cover translation and parsing under substantial compute; architecture ablations are development-set point estimates, not uncertainty-aware general scaling laws or controlled comparisons against recurrent models at equal compute.

\end{document}
